% Reviewed translation: Section 6, PDF pages 372--376 / printed pages 358--362.
\MathBookSourcePage{372}
前面各章已经对此作了详尽研究. 因此, 我们将重点讨论没有不可压约束且纯粹是非线性的
Problem~\eqref{eq:navier-alternating-nonlinear-step}.

Problem~\eqref{eq:navier-alternating-nonlinear-step} 具有如下形式:
\begin{equation}
\label{eq:navier-least-squares-model}
\left\{
\begin{aligned}
-\nu\Delta\boldsymbol u+
\sum_{j=1}^N u_j\partial\boldsymbol u/\partial x_j+c\boldsymbol u
&=\boldsymbol f &&\text{in }\Omega,\\
\boldsymbol u&=\boldsymbol0 &&\text{on }\Gamma.
\end{aligned}
\right.
\end{equation}
其左端同时含有一个线性椭圆项和一个非线性项. 该问题可以方便地推广, 使其适合
Subsection~\ref{subsec:navier-general-descent} 的抽象框架.

以 $\|\cdot\|_*$ 表示熟知的 $X'$ 的对偶范数, 并设
$A\in\mathcal L(X;X')$ 是 $X$ 上的对称 $X$-椭圆算子, 即
\begin{equation}
\label{eq:navier-least-squares-ellipticity}
\langle Av,v\rangle\geq\gamma\|v\|_X^2
\qquad\forall v\in X,\quad\gamma>0.
\end{equation}
设 $G$ 是从 $X$ 到 $X'$ 的 $\mathcal C^p$ 映射($p\geq2$), 并置
\begin{equation}
\label{eq:navier-least-squares-operator}
F(v)=Av+G(v),
\end{equation}
它显然是从 $X$ 到 $X'$ 的 $\mathcal C^p$ 映射. 我们的问题是:
\begin{equation}
\label{eq:navier-least-squares-abstract-problem}
\text{求 }u\in X\text{ 使 }F(u)=0.
\end{equation}
显然, Problem~\eqref{eq:navier-least-squares-model} 是
\eqref{eq:navier-least-squares-abstract-problem} 的特殊情形, 其中
\[
X=H_0^1(\Omega)^N,
\qquad A=-\nu\Delta+c,
\qquad G(\boldsymbol v)=\sum_{j=1}^Nv_j\partial\boldsymbol v/\partial x_j-
\boldsymbol f.
\]

我们打算把 Problem~\eqref{eq:navier-least-squares-abstract-problem} 替换为一个等价的
最小二乘问题来求解. 为此, 考虑泛函
\[
J(v)=\frac12\|F(v)\|_{X'}^2,
\]
其中 $\|\cdot\|_{X'}$ 定义为
\begin{equation}
\label{eq:navier-least-squares-dual-norm}
\|f\|_{X'}=\langle A^{-1}f,f\rangle^{1/2}.
\end{equation}
下面将看到, 由于 $A$ 对称且椭圆, $\|\cdot\|_{X'}$ 是 $X'$ 上与对偶范数等价
的范数, 相应的泛函 $J$ 严格凸且有唯一极小值.

\setcounter{statement}{0}
\begin{Lemma}
\label{lem:navier-least-squares-equivalent-dual-norm}
映射 $f\mapsto\langle A^{-1}f,f\rangle^{1/2}$ 是 $X'$ 上与对偶范数等价的范数.
\end{Lemma}

\begin{Proof}
由 \,\eqref{eq:navier-least-squares-ellipticity} 推得
\[
\langle A^{-1}f,f\rangle\geq\gamma\|A^{-1}f\|_X^2,
\]
即
\[
\|A^{-1}f\|_X\leq\frac1\gamma\|f\|_*.
\]
\MathBookSourcePage{373}
所以一方面
\[
\langle A^{-1}f,f\rangle\leq\frac1\gamma\|f\|_*^2.
\]
另一方面, 置
\[
K=\|A\|_{\mathcal L(X;X')}.
\]
于是
\[
\|f\|_*\leq K\|A^{-1}f\|_X,
\]
从而
\[
\langle A^{-1}f,f\rangle\geq\frac\gamma{K^2}\|f\|_*^2.
\]
综上,
\begin{equation}
\label{eq:navier-least-squares-norm-equivalence}
\frac\gamma{K^2}\|f\|_*^2
\leq\langle A^{-1}f,f\rangle
\leq\frac1\gamma\|f\|_*^2.
\end{equation}
\end{Proof}

\setcounter{statement}{2}
\begin{Theorem}
\label{thm:navier-least-squares-local-convexity}
设 $u$ 是 Problem~\eqref{eq:navier-least-squares-abstract-problem} 的一个非奇异解.
则由
\begin{equation}
\label{eq:navier-least-squares-functional}
J(v)=\frac12\langle A^{-1}(F(v)),F(v)\rangle
\end{equation}
定义的泛函 $J$ 在 $u$ 的一个邻域内严格凸.
\end{Theorem}

\begin{Proof}
考虑到 $A^{-1}$ 的对称性, $J$ 的前两个导数具有如下表达式:
\begin{equation}
\label{eq:navier-least-squares-first-derivative}
\begin{aligned}
DJ(v)\mathbin{\cdot}w
&=\langle A^{-1}(DF(v)\mathbin{\cdot}w),F(v)\rangle\\
&=\langle A^{-1}(F(v)),DF(v)\mathbin{\cdot}w\rangle,
\end{aligned}
\end{equation}
\begin{equation}
\label{eq:navier-least-squares-second-derivative}
D^2J(v)\mathbin{\cdot}(w,z)
=\langle A^{-1}(DF(v)\mathbin{\cdot}z),DF(v)\mathbin{\cdot}w\rangle
+\langle A^{-1}(F(v)),D^2F(v)\mathbin{\cdot}(w,z)\rangle.
\end{equation}
现在回顾, $u$ 是 Problem~\eqref{eq:navier-least-squares-abstract-problem} 的非奇异解,
是指
\[
F(u)=0,
\qquad DF(u)\text{ 是从 }X\text{ 到 }X'\text{ 的同构}.
\]
因此 \,\eqref{eq:navier-least-squares-second-derivative} 给出
\begin{equation}
\label{eq:navier-least-squares-hessian-at-solution}
D^2J(u)\mathbin{\cdot}(w,w)
=\langle A^{-1}(DF(u)\mathbin{\cdot}w),DF(u)\mathbin{\cdot}w\rangle.
\end{equation}
所以 \,\eqref{eq:navier-least-squares-norm-equivalence} 蕴含
\[
\begin{aligned}
D^2J(u)\mathbin{\cdot}(w,w)
&\geq\frac\gamma{K^2}\|DF(u)\mathbin{\cdot}w\|_*^2\\
&\geq\delta\|w\|_X^2,
\qquad \delta>0,
\end{aligned}
\]
因为 $DF(u)$ 是从 $X$ 到 $X'$ 的同构. 但映射 $D^2J$ 在 $X$ 中连续
(因为 $F$ 是 $\mathcal C^p$ 映射); 因而存在 $\rho>0$, 使得对所有
$v\in S(u;\rho)=\{v\in X;\ \|v-u\|_X\leq\rho\}$, 有
\[
\|D^2J(u)-D^2J(v)\|\leq\delta/2.
\]
所以
\MathBookSourcePage{374}
\begin{equation}
\label{eq:navier-least-squares-local-strong-convexity}
D^2J(v)\mathbin{\cdot}(w,w)\geq\frac\delta2\|w\|_X^2
\qquad\forall v\in S(u;\rho),\quad\forall w\in X,
\end{equation}
即 $J$ 在 $S(u;\rho)$ 中严格凸.
\end{Proof}

因此, Problem~\eqref{eq:navier-least-squares-abstract-problem} 等价于求解
\begin{equation}
\label{eq:navier-least-squares-minimization}
\inf_{v\in S(u;\rho)}\frac12\langle A^{-1}F(v),F(v)\rangle,
\end{equation}
而该解可以通过 Subsection~\ref{subsec:navier-general-descent} 的梯度法和共轭梯度法得到.
事实上, 假设 $D^2G$ 在 $X$ 的所有有界子集上有界, 从而 $F$, $DF$ 和 $D^2F$
也在那里有界; 并假设 $u$ 是 \,\eqref{eq:navier-least-squares-abstract-problem}
的一个非奇异解. 那么已经知道, \eqref{eq:navier-descent-strong-convexity} 的第一部分
在球 $S(u;\rho)$ 上成立, 而第二部分由
\eqref{eq:navier-least-squares-hessian-at-solution},
\eqref{eq:navier-least-squares-norm-equivalence}, $DF(u)$ 的同构性质和 $D^2J$
的连续性得到. 此外, $F$ 的有界性蕴含 $J$ 在 $S(u;\rho)$ 中有界. 因而,
选取 $S(u;\rho)$ 中的初值 $u^0$, 并置 $J^0=J(u^0)$, 可以取
\[
D=\{v\in X;\ J(v)\leq J(u^0)\}\cap S(u;\rho),
\]
Theorems~\ref{thm:navier-general-descent-convergence} 和
\ref{thm:navier-conjugate-gradient-convergence} 就保证了梯度算法和共轭梯度算法的收敛性.

下面考察简单梯度法的实际实现. $A$ 的对称性和椭圆性促使我们在 $X$ 上取标量积
\[
(u,v)_X=\langle Au,v\rangle
\]
以及相应的范数 $\|u\|_X=\langle Au,u\rangle^{1/2}$. 因此, 根据
\eqref{eq:navier-least-squares-first-derivative}, 梯度 $g(v)$ 由
\[
\begin{aligned}
\langle Ag(v),w\rangle
&=\langle DJ(v),w\rangle\\
&=\langle A^{-1}F(v),DF(v)\mathbin{\cdot}w\rangle\\
&=\langle(DF(v))'A^{-1}F(v),w\rangle
\end{aligned}
\]
定义, 即
\begin{equation}
\label{eq:navier-least-squares-gradient}
g(v)=A^{-1}(DF(v))'A^{-1}F(v).
\end{equation}
所以, 简单梯度算法的一步可以分解为下列操作:
\begin{enumerate}
\item[1$^\circ$)] 计算
\[
z^m=A^{-1}F(u^m),
\qquad
g^m=A^{-1}(DF(u^m))'z^m;
\]
\item[2$^\circ$)] 关于 $\rho$ 极小化 $J(u^m-\rho g^m)$, 其中
\[
J(u^m-\rho g^m)
=\frac12\langle A^{-1}F(u^m-\rho g^m),F(u^m-\rho g^m)\rangle.
\]
\end{enumerate}
每次迭代需要求解两个与算子 $A$ 有关的线性问题, 并确定 $\rho^m$. 作为例子,
下面明确给出 Problem~\eqref{eq:navier-least-squares-model} 的 $\rho^m$ 的计算.

\MathBookSourcePage{375}
首先注意, 因为 $F$ 是二次多项式, 映射 $\rho\mapsto J(v-\rho w)$ 是四次多项式.
所以 $J(u^m-\rho g^m)$ 的 Taylor 展开简化为
\begin{equation}
\label{eq:navier-least-squares-quartic-expansion}
\begin{aligned}
J(u^m-\rho g^m)
={}&J(u^m)-\rho DJ(u^m)\mathbin{\cdot}g^m
+\frac{\rho^2}{2}D^2J(u^m)\mathbin{\cdot}(g^m,g^m)\\
&-\frac{\rho^3}{6}D^3J(u^m)\mathbin{\cdot}(g^m,g^m,g^m)\\
&+\frac{\rho^4}{24}D^4J(u^m)\mathbin{\cdot}(g^m,g^m,g^m,g^m),
\end{aligned}
\end{equation}
其中 $J$ 的三阶和四阶导数具有简单的表达式
\[
D^3J(u^m)\mathbin{\cdot}(g^m,g^m,g^m)
=3\langle A^{-1}DF(u^m)\mathbin{\cdot}g^m,
D^2F(u^m)\mathbin{\cdot}(g^m,g^m)\rangle,
\]
\[
D^4J(u^m)\mathbin{\cdot}(g^m,g^m,g^m,g^m)
=3\langle A^{-1}D^2F(u^m)\mathbin{\cdot}(g^m,g^m),
D^2F(u^m)\mathbin{\cdot}(g^m,g^m)\rangle.
\]

综上, 为求解 Problem~\eqref{eq:navier-least-squares-model}, 简单梯度算法的每次迭代
按如下步骤进行:
\begin{enumerate}
\item[1$^\circ$)] 给定 $\boldsymbol u^m\in H_0^1(\Omega)^N$, 计算
$\boldsymbol z^m\in H^1(\Omega)^N$:
\[
\left\{
\begin{aligned}
-\nu\Delta\boldsymbol z^m+c\boldsymbol z^m
&=-\nu\Delta\boldsymbol u^m+c\boldsymbol u^m
+\sum_{j=1}^Nu_j^m(\partial\boldsymbol u^m/\partial x_j)-\boldsymbol f
&&\text{in }\Omega,\\
\boldsymbol z^m&=\boldsymbol0&&\text{on }\Gamma;
\end{aligned}
\right.
\]
\item[2$^\circ$)] 求 $\boldsymbol g^m\in H_0^1(\Omega)^N$, 使
\[
\begin{aligned}
\nu(\operatorname{grad}\boldsymbol g^m,\operatorname{grad}\boldsymbol v)
+c(\boldsymbol g^m,\boldsymbol v)
={}&\nu(\operatorname{grad}\boldsymbol z^m,\operatorname{grad}\boldsymbol v)
+c(\boldsymbol z^m,\boldsymbol v)\\
&+\sum_{j=1}^N
(u_j^m\partial\boldsymbol v/\partial x_j
+v_j\partial\boldsymbol u^m/\partial x_j,\boldsymbol z^m)
\end{aligned}
\]
对所有 $\boldsymbol v\in H_0^1(\Omega)^N$ 成立, 且
$\boldsymbol g^m=\boldsymbol0$ on $\Gamma$;
\item[3$^\circ$)] 计算
\begin{equation}
\label{eq:navier-least-squares-functional-value}
J(\boldsymbol u^m)=\frac12\{\nu|\boldsymbol z^m|_{1,\Omega}^2
+c\|\boldsymbol z^m\|_{0,\Omega}^2\},
\end{equation}
\begin{equation}
\label{eq:navier-least-squares-gradient-value}
DJ(\boldsymbol u^m)\mathbin{\cdot}\boldsymbol g^m
=\frac12\{\nu|\boldsymbol g^m|_{1,\Omega}^2
+c\|\boldsymbol g^m\|_{0,\Omega}^2\};
\end{equation}
\item[4$^\circ$)] 求 $\boldsymbol v^m\in H^1(\Omega)^N$:
\[
\left\{
\begin{aligned}
-\nu\Delta\boldsymbol v^m+c\boldsymbol v^m
&=-\nu\Delta\boldsymbol g^m+c\boldsymbol g^m
+\sum_{j=1}^N
(u_j^m\partial\boldsymbol g^m/\partial x_j
+g_j^m\partial\boldsymbol u^m/\partial x_j)
&&\text{in }\Omega,\\
\boldsymbol v^m&=\boldsymbol0&&\text{on }\Gamma;
\end{aligned}
\right.
\]
\item[5$^\circ$)] 计算
\[
\boldsymbol t^m=D^2F(\boldsymbol u^m)\mathbin{\cdot}
(\boldsymbol g^m,\boldsymbol g^m)
=2\sum_{j=1}^Ng_j^m\partial\boldsymbol g^m/\partial x_j,
\]
\begin{equation}
\label{eq:navier-least-squares-second-derivative-value}
D^2J(\boldsymbol u^m)\mathbin{\cdot}(\boldsymbol g^m,\boldsymbol g^m)
=\nu|\boldsymbol v^m|_{1,\Omega}^2
+c\|\boldsymbol v^m\|_{0,\Omega}^2+(\boldsymbol z^m,\boldsymbol t^m),
\end{equation}
\begin{equation}
\label{eq:navier-least-squares-third-derivative-value}
D^3J(\boldsymbol u^m)\mathbin{\cdot}(\boldsymbol g^m,\boldsymbol g^m,\boldsymbol g^m)
=3(\boldsymbol t^m,\boldsymbol v^m);
\end{equation}
\end{enumerate}

\MathBookSourcePage{376}
\begin{enumerate}
\item[6$^\circ$)] 求 $\boldsymbol w^m\in H^1(\Omega)^N$:
\[
\left\{
\begin{aligned}
-\nu\Delta\boldsymbol w^m+c\boldsymbol w^m&=\boldsymbol t^m&&\text{in }\Omega,\\
\boldsymbol w^m&=\boldsymbol0&&\text{on }\Gamma;
\end{aligned}
\right.
\]
\item[7$^\circ$)] 计算
\begin{equation}
\label{eq:navier-least-squares-fourth-derivative-value}
D^4J(\boldsymbol u^m)\mathbin{\cdot}
(\boldsymbol g^m,\boldsymbol g^m,\boldsymbol g^m,\boldsymbol g^m)
=\nu|\boldsymbol w^m|_{1,\Omega}^2
+c\|\boldsymbol w^m\|_{0,\Omega}^2;
\end{equation}
\item[8$^\circ$)] 求方程
\[
\frac{dJ(\boldsymbol u^m-\rho^m\boldsymbol g^m)}{d\rho}=0
\]
的正根 $\rho^m$, 并如下更新 $\boldsymbol u^m$:
\[
\boldsymbol u^{m+1}=\boldsymbol u^m-\rho^m\boldsymbol g^m.
\]
\end{enumerate}
每次迭代需要求解四个 Dirichlet 问题.

共轭梯度算法的实现与上述过程非常相似, 留作 Exercise.

\subsection{Newton 方法与延拓法}
\label{subsec:navier-newton-continuation}

这里讨论的方法旨在求解完整的 Navier--Stokes 方程, 即同时处理不可压性和非线性.
更一般地, 我们要解 Subsection~\ref{subsec:navier-nonsingular-abstract} 中引入的一类方程,
即
\begin{equation}
\label{eq:navier-newton-parameterized-problem}
F(\lambda,u)=0,
\end{equation}
其中 $F$ 是定义在 $\Lambda\times X$ 上, 取值于 $\mathscr X$ 的 $\mathcal C^p$ 映射
($p\geq1$), $X$ 和 $\mathscr X$ 是两个 Banach 空间, $\Lambda$ 是 $\mathbb R$
中的一个区间. 暂时固定 $\lambda$, 并假设 $u=u(\lambda)\in X$ 是
\eqref{eq:navier-newton-parameterized-problem} 的一个非奇异解, 即
\[
F(\lambda,u)=0,
\qquad D_uF(\lambda,u)\text{ 是从 }X\text{ 到 }\mathscr X\text{ 的同构}.
\]
由逆函数 Theorem(也由 Subsection~\ref{subsec:navier-nonsingular-branch-approximation}
的材料)可知, 存在一个闭球 $S(u;\alpha)$, 方程
\eqref{eq:navier-newton-parameterized-problem} 在其中除 $u$ 外没有其他解.

由于 $u$ 是 \,\eqref{eq:navier-newton-parameterized-problem} 的孤立解, 且 $F$
至少可微, 逼近 $u$ 的一种有效方法是 Newton 算法: 从初始猜测 $u^0$ 出发,
由下式在 $X$ 中构造序列 $(u^n)$:
\begin{equation}
\label{eq:navier-newton-iteration}
u^{n+1}=u^n-[D_uF(\lambda,u^n)]^{-1}\mathbin{\cdot}F(\lambda,u^n),
\qquad n\geq0,
\end{equation}
或等价地
\[
D_uF(\lambda,u^n)\mathbin{\cdot}(u^{n+1}-u^n)=-F(\lambda,u^n).
\]
由于 $D_uF(\lambda,u)$ 是线性算子, Newton 方法的每一步都需要求解一个与
$D_uF(\lambda,u^n)$ 有关的不同线性问题. 如果代价过高,
